\chapter{INTRODUCTION}
\thispagestyle{empty}

\section{Portee du projet}

Ce projet etudie la recommandation d'article a article sur un graphe Wikipedia
personnalise consacre a l'informatique. L'objectif n'est pas une recommandation
personnalisee fondee sur le comportement d'utilisateurs, mais une recuperation
guidee par requete a partir d'un article source connu: etant donne un article
Wikipedia du domaine de l'informatique, le systeme doit proposer une liste
classee d'articles lies qui soient a la fois semantiquement pertinents et bien
connectes structurellement dans le graphe.

Le travail s'appuie sur le benchmark Wiki-CS~\cite{mernyei2020wikics}, mais
l'etend de maniere importante. Au lieu de traiter les noeuds comme de simples
objets de graphe, le projet les rehydrate avec le texte reel des articles
Wikipedia, calcule des embeddings semantiques a partir de ce texte, apprend des
representations structurelles a partir des hyperliens et evalue la maniere dont
ces signaux peuvent etre combines pour la recommandation. Le projet comprend
egalement un prototype web local qui expose les modeles de recuperation via une
interface de recherche et de previsualisation.

\section{Motivation}

Wikipedia contient une tres grande quantite de connaissances ouvertes, mais sa
navigation reste souvent peu efficace pour l'apprentissage. Les moteurs de
recherche peuvent trouver une page correspondant a un mot-cle, sans pour autant
orienter l'utilisateur vers des prerequis proches, des explications
alternatives ou des sujets relies structurellement. Dans un graphe tel que
Wikipedia, les hyperliens encodent deja une grande quantite d'information
editoriale et conceptuelle. En meme temps, le texte des articles porte une
similarite semantique qui n'apparait pas toujours a travers la seule structure
des liens.

Le projet etudie donc si la combinaison de la structure du graphe et de la
representation semantique du texte peut produire de meilleures recommandations
d'articles que chacune de ces sources d'information prise isolement. La
motivation pedagogique plus large est de passer d'une recuperation isolee
d'articles a une exploration coherente d'un domaine de connaissances.

\section{De l'idee initiale a la direction finale}

L'orientation initiale du projet etait plus large que le systeme finalement
mis en oeuvre. Elle comprenait l'idee d'un generateur complet de parcours
d'apprentissage combinant reseaux de neurones sur graphes, recuperation
semantique et construction hierarchique fondee sur un LLM. Au cours du
developpement, le centre de gravite pratique s'est deplace. Le travail realise
s'est revele le plus solide dans trois directions:
\begin{itemize}
    \item la construction d'un jeu de donnees de graphe enrichi et reutilisable;
    \item la comparaison empirique de modeles de recommandation et de prediction
    de liens;
    \item le developpement d'une interface prototype locale pour la recuperation
    d'articles.
\end{itemize}

En consequence, le rapport final presente avant tout le projet comme une etude
empirique de recommandation sur graphe, accompagnee d'une couche optionnelle
d'organisation en parcours d'apprentissage, plutot que comme une plateforme
utilisateur completement aboutie.

\section{Contributions principales}

Le projet apporte les contributions concretes suivantes:
\begin{itemize}
    \item la creation de l'espace experimental Custom-WikiCS, qui etend Wiki-CS
    avec le texte des articles, des caches d'embeddings et des artefacts
    d'evaluation;
    \item l'analyse du graphe nettoye d'articles d'informatique, incluant le
    comportement des degres, l'assortativite, la centralite et la structure
    communautaire;
    \item la construction et l'evaluation de modeles de recommandation
    semantiques, structurels, hybrides et fondes sur des GNN;
    \item la conception d'une couche optionnelle d'organisation en parcours
    d'apprentissage, fondee sur un LLM contraint, une validation cote
    application et un repli heuristique;
    \item l'implementation d'une application web locale pour la recherche par
    titre, la selection de modele, les recommandations classees, la
    previsualisation d'article et l'organisation optionnelle en parcours
    d'apprentissage.
\end{itemize}

\section{Portee finale implementee}

Dans sa forme finale, l'etude s'appuie sur un graphe nettoye comprenant 11,367
noeuds actifs et 291,039 aretes orientees. Les principales familles de modeles
sont:
\begin{itemize}
    \item les embeddings semantiques BGE-M3~\cite{chen2025bgem3};
    \item les embeddings structurels de type Node2Vec~\cite{grover2016node2vec};
    \item les baselines structurelles GCN et GraphSAGE~\cite{kipf2017semi,hamilton2017inductive};
    \item les variantes hybrides de recuperation semantico-structurelle et de
    fusion de rang.
\end{itemize}

Les Graph Attention Networks (GAT)~\cite{velickovic2018graph} restent
conceptuellement pertinents, mais ne constituent pas la contribution mise en
oeuvre principale de cette etude. Les resultats les plus solides et les plus
stables proviennent du cadre revise.

\section{Organisation du rapport}

La suite du rapport est organisee comme suit. Le Chapitre 2 presente la
litterature principale relative a la recuperation semantique, a
l'apprentissage de representations sur graphe et a la recommandation sur
Wikipedia. Le Chapitre 3 definit le probleme de recommandation et les
contraintes du projet. Le Chapitre 4 presente la methodologie, l'architecture
du systeme, les familles de modeles et la conception actuelle de l'application.
Le Chapitre 5 resume le dispositif experimental et les principaux resultats
quantitatifs. Le Chapitre 6 discute l'interpretation de ces resultats, y
compris le passage du premier dispositif d'evaluation vers le protocole corrige
et l'evolution de l'organisateur fonde sur LLM. Le Chapitre 7 conclut le
rapport et ouvre sur les perspectives futures.
